\documentclass[11pt]{amsart}

\usepackage[T1]{fontenc}
\usepackage{lmodern}
\usepackage{microtype}
\usepackage{amsmath,amssymb,amsthm}
\usepackage{array}
\usepackage[hidelinks]{hyperref}

\newtheorem{theorem}{Theorem}
\newtheorem{proposition}[theorem]{Proposition}
\newtheorem{corollary}[theorem]{Corollary}

\title[A tree refuting a $GRM_{-2}$ lower bound]{A 12-Vertex Tree Refuting the
\texorpdfstring{$GRM_{-2}$}{GRM(-2)} Lower Bound for Trees of Maximum Degree Four}
\date{}
\subjclass[2020]{Primary 05C09; Secondary 05C05, 05C07, 05C35, 92E10}
\keywords{general reduced second Zagreb index, $GRM_\lambda$, reduced Zagreb index,
molecular tree, maximum degree, extremal tree, counterexample}

\hypersetup{
  pdftitle={A 12-vertex tree refuting the GRM(-2) lower bound for trees of maximum degree four}
}

\begin{document}

\begin{abstract}
For a graph $G$ and $\lambda\in\mathbb{R}$, the general reduced second Zagreb index is
$GRM_\lambda(G)=\sum_{uv\in E(G)}(\deg u+\lambda)(\deg v+\lambda)$. Ba\v{s}i\'{c} and
Ili\'{c} assert that every tree $T$ of order $n$ with $\Delta(T)=4$ and
$n\not\equiv\{1,2,3\}\pmod 4$ satisfies $GRM_{-2}(T)\ge -n$, and that this value is attained.
We exhibit a tree on $12$ vertices with maximum degree $4$ and $GRM_{-2}=-13<-12$, so the
assertion is false, and we exhibit an infinite family showing
that for every $n\equiv 0\pmod 4$ with $n\ge 16$ the assertion fails by $2$. Consequently the
claimed extremal family $TT^4_{opt}(k)$ is not extremal for any $k\ge 2$. The refutation is a
finite hand computation on eleven edges. We do not determine the exact minimum at any order,
we do not decide whether the assertion holds at $n=8$, and
we make no claim about the source's other results.
\end{abstract}

\maketitle

\section{The claim under test, and its locator}\label{sec:claim}

For a graph $G$ and a real parameter $\lambda$, Horoldagva, Buyantogtokh, Das and Lee
\cite{Horoldagva} introduced the \emph{general reduced second Zagreb index}
\begin{equation}\label{eq:def}
  GRM_\lambda(G)=\sum_{uv\in E(G)}\bigl(\deg(u)+\lambda\bigr)\bigl(\deg(v)+\lambda\bigr),
\end{equation}
so that
\begin{equation}\label{eq:grm2}
  GRM_{-2}(T)=\sum_{uv\in E(T)}\bigl(\deg(u)-2\bigr)\bigl(\deg(v)-2\bigr).
\end{equation}
Write $m_{ij}$ for the number of edges of $T$ whose endpoint degrees are $i$ and $j$, and
$n_i$ for the number of vertices of degree $i$.

Ba\v{s}i\'{c} and Ili\'{c} \cite{BasicIlic} devote their Section~4 to the minimum of
$GRM_{-2}$ over trees of maximum degree $4$. All references below are to line numbers of the
arXiv \LaTeX{} e-print of \cite{BasicIlic}, version~1; the class is fixed at line~797 by the
phrase ``\emph{for any tree $T$ with a maximal degree of $\Delta=4$}'', and the section
heading at line~766 reads ``\emph{Minimum value of $GRM_{-2}$ in the class of molecular trees
with maximal degree 4}''.
The source's Section~4 rests on its own reduction, the displayed equation at line~800,
\begin{equation}\label{eq:main2}
  -GRM_{-2}(T)=-3m_{12}-m_{13}-m_{22}-m_{34}-3m_{44}+n-n_3+3 ,
\end{equation}
which in turn comes from its pre-substitution formula at line~795,
\begin{equation}\label{eq:l795}
  -GRM_{-2}(T)=m_{13}+2m_{14}-m_{33}-2m_{34}-4m_{44}.
\end{equation}
The section closes four residue cases in turn. The last of them, at line~838, is quoted here
character for character:
\begin{verbatim}
Therefore, based on equation (\ref{main2}), if $n
\not\equiv \{1, 2, 3\} \pmod{4}$, it follows that
$GRM_{-2}(T) \geq -n$.
\end{verbatim}
With the class of line~797, this is the following statement, which we refute.

\begin{quote}
\textbf{Claim} (\cite{BasicIlic}, line~838). \emph{For every tree $T$ of order $n$ with
$\Delta(T)=4$ and $n\equiv 0\pmod 4$ one has $GRM_{-2}(T)\ge -n$.}
\end{quote}

Line~842 adds that the value $-n=-(4k+4)$ \emph{is attained}, at parameters listed at
lines~843 and~844, and line~845 names the family $TT^4_{opt}(k)$ of trees realising it.
The adjacent improvement of Dehgardi and Klav\v{z}ar
\cite{DehgardiKlavzar}, reference \textup{[3]} of \cite{BasicIlic}, is restricted to
$\lambda\ge-\tfrac12$ and so says nothing about $\lambda=-2$.

\section{The counterexample}\label{sec:ce}

\begin{theorem}\label{thm:main}
There is a tree $W_1$ on $n=12$ vertices with $\Delta(W_1)=4$ and $GRM_{-2}(W_1)=-13$. Since
$12\equiv 0\pmod 4$ and $-13<-12=-n$, the Claim of Section~\textup{\ref{sec:claim}} is false.
\end{theorem}

\begin{proof}
Let $W_1$ have vertex set $\{0,1,\dots,11\}$ and the eleven edges
\begin{equation}\label{eq:w1}
  01,\quad 02,\quad 03,\quad 04,\quad 15,\quad 56,\quad 57,\quad
  68,\quad 89,\quad 8\,\text{-}\,10,\quad 8\,\text{-}\,11 .
\end{equation}
Pictorially, $0$ and $8$ are the two branch vertices, joined by the path $0,1,5,6,8$, with a
single extra leaf at $5$:
\[
\begin{array}{ccccccccc}
 2\ \ 3\ \ 4 & & & & & & & & 9\ \ 10\ \ 11\\[-2pt]
 \diagdown\ |\ \diagup & & & & & & & & \diagdown\ \ |\ \ \diagup\\[-2pt]
 0 & \text{---} & 1 & \text{---} & 5 & \text{---} & 6 & \text{---} & 8\\[-2pt]
 & & & & |\\[-2pt]
 & & & & 7
\end{array}
\]

\emph{$W_1$ is a tree.} It has $12$ vertices and $11=n-1$ edges, no loop and no repeated edge
in \eqref{eq:w1}, and every vertex is reached from $0$ by
$0\to\{1,2,3,4\}$, $1\to 5$, $5\to\{6,7\}$, $6\to 8$, $8\to\{9,10,11\}$. A connected graph
with $n-1$ edges is a tree.

\emph{Degrees.} From \eqref{eq:w1}, $\deg 0=\deg 8=4$, $\deg 5=3$, $\deg 1=\deg 6=2$, and
$2,3,4,7,9,10,11$ are leaves. The degree sequence is $(4,4,3,2,2,1^7)$, whose sum is
$22=2\cdot 11$. In particular $\Delta(W_1)=4$ exactly, attained at $0$ and at $8$, so $W_1$
lies in the class of line~797 under either reading of it, $\Delta=4$ or $\Delta\le 4$.

\emph{The value.} Evaluate \eqref{eq:grm2} edge by edge. Only the seven edges incident with a
leaf contribute, because an edge with an endpoint of degree $2$ contributes $0$.
\[
\begin{array}{l|c|r|r}
 \text{edges} & \text{end degrees} & \text{each} & \text{subtotal}\\\hline
 02,\ 03,\ 04 & (4,1) & (4-2)(1-2)=-2 & -6\\
 89,\ 8\text{-}10,\ 8\text{-}11 & (4,1) & -2 & -6\\
 57 & (3,1) & (3-2)(1-2)=-1 & -1\\
 01 & (4,2) & (4-2)(2-2)=0 & 0\\
 15 & (2,3) & 0 & 0\\
 56 & (3,2) & 0 & 0\\
 68 & (2,4) & 0 & 0\\\hline
 \multicolumn{3}{r|}{\text{total}} & -13
\end{array}
\]
Hence $GRM_{-2}(W_1)=-13$, while the Claim asserts $GRM_{-2}(W_1)\ge -12$.
\end{proof}

Two remarks fix the scope of Theorem~\ref{thm:main}. First, $W_1$ itself violates only the
Claim of line~838: the source's earlier conclusions at line~822 ($GRM_{-2}\ge-(n+2)$ for
$n\ne 4k+1$) and line~830 ($GRM_{-2}\ge-(n+1)$ for $n\not\equiv\{1,2\}\pmod 4$) both survive
at $W_1$, the second with equality, since $-13=-(n+1)$. The tree $W_2$ of
Section~\ref{sec:family} violates line~830 as well. Second, at $n=8$ the bound $-n$ is
attained, by the double star obtained by joining two vertices of degree $4$, which has
$m_{14}=6$, $m_{44}=1$ and $GRM_{-2}=6\cdot(-2)+1\cdot 4=-8=-n$; we do not decide whether
$GRM_{-2}(T)\ge-8$ holds for every tree $T$ of order $8$ with $\Delta(T)=4$, and nothing here
depends on that.

The edge partition of $W_1$ is $m_{13}=1$, $m_{14}=6$, $m_{23}=2$, $m_{24}=2$, all other
$m_{ij}=0$, with $n_1=7$, $n_2=2$, $n_3=1$, $n_4=2$; both \eqref{eq:main2} and
\eqref{eq:l795} return $-GRM_{-2}(W_1)=13$.

\section{An infinite family}\label{sec:family}

\begin{proposition}\label{prop:family}
For every $k\ge 3$ there is a tree $F(k)$ of order $n=4k+4$ with $\Delta=4$ and
$GRM_{-2}(F(k))=-(n+2)$.
\end{proposition}

\begin{proof}
Take any tree $S$ (a \emph{skeleton}) on $k+1$ nodes in which a distinguished node $v$ has
degree $3$ and every node has degree at most $4$; such an $S$ exists for every $k\ge 3$, for
instance the tree on $\{0,1,\dots,k\}$ with edges $01,02,03$ together with the path edges
$i(i+1)$ for $3\le i\le k-1$ (no path edge at all when $k=3$, where $S=K_{1,3}$), whose
node $v=0$ has degree $3$. Subdivide every edge of $S$ exactly once, which creates $k$ new
vertices of degree $2$, and then attach $4-d_S(u)$ pendant vertices to each node $u\ne v$,
where $d_S(u)$ is the degree of $u$ in $S$.

In the resulting tree $F(k)$ the node $v$ has degree $3$ and all three of its neighbours are
subdivision vertices, so $n_3=1$ and $m_{13}=0$; each node $u\ne v$ has degree $4$, so
$n_4=k$; the $k$ subdivision vertices give $n_2=k$; and every edge of $F(k)$ either joins a
degree-$2$ vertex to something or joins a degree-$4$ vertex to a leaf, whence
$m_{12}=m_{22}=m_{33}=m_{34}=m_{44}=0$. The skeleton degree sum is $2k$, so
$\sum_{u\ne v}d_S(u)=2k-3$ and the number of pendant vertices is
$n_1=m_{14}=4k-(2k-3)=2k+3$. Therefore
\[
  n=n_1+n_2+n_3+n_4=(2k+3)+k+1+k=4k+4 ,
\]
and by \eqref{eq:grm2} only the $m_{14}$ edges contribute, each $-2$, so
$GRM_{-2}(F(k))=-2(2k+3)=-(4k+6)=-(n+2)$.
\end{proof}

For $k=3$ this is the tree $W_2$ on $16$ vertices with edges
\[
 01,\ 02,\ 03,\ 04,\ 15,\ 56,\ 57,\ 68,\ 79,\ 8\text{-}10,\ 8\text{-}11,\ 8\text{-}14,\
 9\text{-}12,\ 9\text{-}13,\ 9\text{-}15,
\]
with $\Delta=4$ attained at $0,8,9$, $m_{14}=9$ and $GRM_{-2}(W_2)=-18<-16$. Note that $W_2$
violates the source's line~830 as well, since $-18<-17=-(n+1)$, while $W_1$ does not.

\begin{corollary}
The family $TT^4_{opt}(k)$ of \textup{\cite{BasicIlic}}, line~845, is not the set of extremal
trees of order $n=4k+4$ for any $k\ge 2$.
\end{corollary}

\begin{proof}
Both parameter sets the source assigns to its optimal tree have $m_{13}=m_{33}=m_{34}=0$; the
first (line~843) has $m_{14}=2k+2$ and $m_{44}=0$, the second (line~844) has $m_{14}=2k+4$ and
$m_{44}=1$. By \eqref{eq:l795} they give $-GRM_{-2}=2(2k+2)=4k+4=n$ and
$-GRM_{-2}=2(2k+4)-4=4k+4=n$ respectively, i.e. $GRM_{-2}=-n$ in both cases. By
Theorem~\ref{thm:main} and Proposition~\ref{prop:family} the minimum at $n=4k+4$ is at most
$-(n+1)$ for $k=2$ and at most $-(n+2)$ for $k\ge 3$, strictly below $-n$.
\end{proof}

\section{What is not settled}\label{sec:open}

\begin{itemize}
\item \emph{No exact minimum is claimed.} Theorem~\ref{thm:main} and
      Proposition~\ref{prop:family} give upper bounds on
      $\min\{GRM_{-2}(T):|V(T)|=n,\ \Delta(T)=4\}$ only. We do not prove a matching lower
      bound at any order, we do not claim that the trees exhibited here are minimisers, and we
      make no claim about the number of minimisers.
\item \emph{No corrected replacement for $TT^4_{opt}(k)$ is proved here.} The inductive
      structural characterisation at line~846 of \cite{BasicIlic} is not repaired.
\item \emph{The disputed statements are line~838, the line~830 bound at $n\equiv0\pmod4$ with
      $n\ge16$, and the attainment claim at lines~842--846.} We assert nothing about the
      source's $\Delta=3$ results, its $\lambda\ge-1$ material, or its $\Delta=4$ conclusions
      at $n\equiv 1,2,3\pmod 4$.
\item \emph{``Molecular tree'' is never separately defined} in \cite{BasicIlic}; we read the
      class as it is fixed at line~797. All trees exhibited here have $\Delta$ exactly $4$, so
      the refutation is insensitive to reading the class as $\Delta\le 4$ instead.
\item \emph{What the accompanying program does.} The recorded run in the folder of this note
      recomputes $GRM_{-2}$ and the degree and edge counts of $W_1$, $W_2$, the tree $F(4)$ on
      $20$ vertices and the $n=8$ double star from their edge lists, and instantiates $F(k)$
      for $k=3,\dots,12$. It enumerates no trees, so it decides no minimum and does not
      address $n=8$ beyond the double star; the general statement of
      Proposition~\ref{prop:family} for all $k\ge 3$ is established by the proof above, not by
      that run.
\item \emph{This corrects a preprint.} \cite{BasicIlic} has one version, no journal reference
      and no erratum at the time of writing. Should a refereed version appear with a repaired
      Section~4, the target of this note would cease to exist in the literature, though nothing
      about $W_1$ would change.
\end{itemize}

\begin{thebibliography}{9}

\bibitem{BasicIlic}
M.~Ba\v{s}i\'{c} and A.~Ili\'{c},
\emph{Further results on the lower bound on reduced Zagreb index of trees},
\href{https://arxiv.org/abs/2604.06044}{arXiv:2604.06044}, 2026.
(All line numbers cited here refer to \texttt{main.tex} of the version-1 e-print,
$966$ lines, $61{,}112$ bytes.)

\bibitem{DehgardiKlavzar}
N.~Dehgardi and S.~Klav\v{z}ar,
\emph{Improved lower bounds on the general reduced second Zagreb index of trees},
preprint, 2023; cited as reference \textup{[3]} of \cite{BasicIlic}. Published form:
N.~Dehgardi, S.~Klav\v{z}ar and M.~Azari,
J. Comb. Optim. \textbf{51} (2026), Paper No.~50
(\href{https://arxiv.org/abs/2605.30402}{arXiv:2605.30402}); it treats
$\lambda\ge-\tfrac12$ only.

\bibitem{Horoldagva}
B.~Horoldagva, L.~Buyantogtokh, K.~C.~Das and S.-G.~Lee,
\emph{On general reduced second Zagreb index of graphs},
Hacettepe J. Math. Stat. \textbf{48} (2019), 1046--1056.

\end{thebibliography}

\end{document}
